\section*{Zusammenfassung}
\addcontentsline{toc}{section}{Zusammenfassung}

Das Dokument entwickelt ein verbessertes Verfahren zur Bundestagswahl aus drei Bausteinen:

\begin{enumerate}
  \item \textbf{Iteratives Verhältniswahlrecht} --- Aufeinanderfolgende Wahlgänge stellen sicher,
    dass es immer einen Wahlgewinner gibt, der mindestens von der Hälfte des Volkes gewählt wurde.
    Proportionalität wird auf die Opposition beschränkt. Koalitionen sind nach wie vor möglich,
    aber auch ohne sie ist Regierungsfähigkeit sichergestellt.

  \item \textbf{Direkte Wahl der Spitzenkandidaten} --- Jede Partei tritt mit drei
    Spitzenkandidaten an. Jeder Wähler gibt eine \enquote{Parteienstimme} (bestimmt die
    Sitzverteilung im Parlament) und eine \enquote{Spitzenkandidatenstimme} ab. Mit der
    Spitzenkandidatenstimme wählt er den Spitzenkandidaten seiner Partei.

  \item \textbf{Personelle Vertretung jedes Wahlkreises} --- Wähler können nicht mehr direkt
    bestimmen, welche Partei ihren Wahlkreis vertreten soll. Das wird durch einen Algorithmus
    festgelegt, der anhand der Parteienstimmenverteilung die relative Wichtigkeit des Wahlkreises
    für eine Partei ermittelt und die Wahlkreise so Parteien zuweist, dass die relative
    Gesamtwichtigkeit maximiert wird.
    Jede Partei muss aber auf Wahlkreisebene drei Abgeordnetenkandidaten stellen. Die Wähler
    können für jede Partei in ihrem Wahlkreis (nicht nur für die eigene), den ihnen genehm\-sten
    Abgeordneten bestimmen. Der Abgeordnetenkandidat der zugewiesenen Partei mit den meisten
    Stimmen gewinnt den Wahlkreis.
\end{enumerate}

Eine Open-Source-Simulation des Verfahrens steht als
\href{https://github.com/danielschwalbe/ipres}{GitHub-Projekt} zur Verfügung.
